\section{Powder-Synthesis Routes}

\subsection{Solid-state reaction and precursor control}

In this framework, ``solid state'' names a powder-forming family, not a complete experiment. The audited set contains an early paper explicitly titled around synthesis of cubic LLZO for solid-state batteries \cite{tan2012synthesis}; because its local record is title-only, it anchors the existence of that route but not a particular recipe or outcome. A more information-rich abstract follows Al-containing LLZO through solid-state processing and reports that lithium inventory and powder surface-to-volume geometry are consequential within that study \cite{heywood2023tailoring}. The contrast between those evidence levels is deliberate: the former locates a topic, while the latter supports an abstract-reported association.

Precursor identity can be varied without changing the route label. One stored abstract compares monoclinic and tetragonal zirconia starting powders under a reported common synthesis path and includes phase and transport readouts \cite{rahmawati2021a}. That record supports a controlled-within-study precursor comparison, but it does not make zirconia polymorph the only variable needed to compare against other studies. Likewise, in situ neutron diffraction is reported as tracking LLZO formation and identifying precursor-state and cooling-history features \cite{rao2015in}. The latter is valuable because it treats phase formation as an observed sequence; any generalization beyond the studied composition and schedule would require its full methods and results.

Accordingly, the solid-state branch should be read through E1 fields: reagent identity and condition, lithium excess, mixing intensity, intermediate calcination, powder geometry, and provenance. If these fields differ or are unavailable, two ``solid-state'' specimens are parallel qualitative evidence, not direct competitors.

\subsection{Solution-mediated and chelation routes}

The collection contains title-level records for nitrogen-free sol--gel preparation of Al-substituted cubic LLZO \cite{rosenkiewitz2015nitrogen} and for sol--gel synthesis coupled to lithium-ion-conductivity measurement \cite{kokal2011sol}. Those titles establish solution-mediated coverage but do not license claims about lower temperature, finer powder, or superior transport. Such outcomes are conditional on the checked evidence.

Three abstract-bearing records permit more specific, still bounded statements. A solution-derived powder study reports subsequent pressure-assisted densification, showing that powder synthesis and pellet processing cannot be credited to one route label alone \cite{sakamoto2013synthesis}. A comparative study places solid-state and sol--gel preparation, with and without Al, in one thermal, structural, and electrochemical analysis \cite{kosir2022comparative}. A PVP-chelated solution route reports predominantly cubic LLZO together with a secondary phase and an impedance-derived conductivity, so ``phase obtained'' should not be shortened to ``phase pure'' \cite{bae2019use}. These records show why a solution route must be carried forward with its later densification and measurement record.

\subsection{Coprecipitation, combustion, and rapid formation}

A coprecipitation paper and a low-temperature powder-processing paper are title-only in the local corpus \cite{hamao2015synthesis,padarti2018low}. A microwave-assisted solution-combustion paper is also title-only and is therefore used only to mark rapid powder formation for Ga-containing LLZO \cite{dandeneau2023rapid}. No temperature, phase fraction, particle-size advantage, or scale claim is assigned to those three records here.

By contrast, the stored abstract for a nonaqueous polymer-combustion study reports undoped and Ta-containing LLZO nanopowders and follows them into later sintering and transport characterization \cite{weller2019nonaqueous}. It therefore supplies an example of a route-to-readout chain rather than proof that combustion intrinsically outperforms a different route. The comparison lesson across P1.1--P1.3 is consistent: powder formation may alter mixing length scale, intermediate chemistry, or particle state, but a defensible route comparison must also preserve composition, thermal budget, green-body preparation, density, and transport definition. Where those common fields are absent, this review leaves the evidence parallel.

\begin{table}[t]
\centering
\caption{Evidence-matrix examples. The final column caps the inference at the locally checked provenance.}
\label{tab:evidence-matrix}
\small
\begin{tabular}{P{0.16\linewidth}P{0.23\linewidth}P{0.22\linewidth}P{0.28\linewidth}}
\toprule
Primary tag & Intervention represented & Coupled readout & Licensed use here \\
\midrule
P1.1 & Zirconia precursor identity & Phase and transport & Within-study precursor comparison, abstract level \cite{rahmawati2021a} \\
P1.2 & Solid-state versus sol--gel, with composition branch & Thermal, structural, electrochemical & Multi-factor comparative case, abstract level \cite{kosir2022comparative} \\
P1.3 & Nonaqueous polymer combustion, undoped/Ta branch & Powder, sintering, transport & Route-to-readout case, abstract level \cite{weller2019nonaqueous} \\
P3.2 & Hot-pressing temperature at fixed pressure & Density and transport partition & Within-study densification case, abstract level \cite{david2015microstructure} \\
M3 & Crucible choice and air exposure & Composition, density, transport & Boundary-condition comparison, abstract level \cite{xia2016ionic} \\
\bottomrule
\end{tabular}
\end{table}
